% Options for packages loaded elsewhere
\PassOptionsToPackage{unicode}{hyperref}
\PassOptionsToPackage{hyphens}{url}
\PassOptionsToPackage{dvipsnames,svgnames,x11names}{xcolor}
%
\documentclass[
  11pt,
]{article}
\usepackage{amsmath,amssymb}
\usepackage{iftex}
\ifPDFTeX
  \usepackage[T1]{fontenc}
  \usepackage[utf8]{inputenc}
  \usepackage{textcomp} % provide euro and other symbols
\else % if luatex or xetex
  \usepackage{unicode-math} % this also loads fontspec
  \defaultfontfeatures{Scale=MatchLowercase}
  \defaultfontfeatures[\rmfamily]{Ligatures=TeX,Scale=1}
\fi
\usepackage{lmodern}
\ifPDFTeX\else
  % xetex/luatex font selection
\fi
% Use upquote if available, for straight quotes in verbatim environments
\IfFileExists{upquote.sty}{\usepackage{upquote}}{}
\IfFileExists{microtype.sty}{% use microtype if available
  \usepackage[]{microtype}
  \UseMicrotypeSet[protrusion]{basicmath} % disable protrusion for tt fonts
}{}
\makeatletter
\@ifundefined{KOMAClassName}{% if non-KOMA class
  \IfFileExists{parskip.sty}{%
    \usepackage{parskip}
  }{% else
    \setlength{\parindent}{0pt}
    \setlength{\parskip}{6pt plus 2pt minus 1pt}}
}{% if KOMA class
  \KOMAoptions{parskip=half}}
\makeatother
\usepackage{xcolor}
\usepackage[margin=1in]{geometry}
\setlength{\emergencystretch}{3em} % prevent overfull lines
\providecommand{\tightlist}{%
  \setlength{\itemsep}{0pt}\setlength{\parskip}{0pt}}
\setcounter{secnumdepth}{-\maxdimen} % remove section numbering
\ifLuaTeX
  \usepackage{selnolig}  % disable illegal ligatures
\fi
\usepackage{bookmark}
\IfFileExists{xurl.sty}{\usepackage{xurl}}{} % add URL line breaks if available
\urlstyle{same}
\hypersetup{
  pdftitle={Prescribed equality in off-grid extensions of G449},
  pdfauthor={Aafi Khan Malek (Independent Researcher, Ahmedabad, India)},
  colorlinks=true,
  linkcolor={Maroon},
  filecolor={Maroon},
  citecolor={Blue},
  urlcolor={Blue},
  pdfcreator={LaTeX via pandoc}}

\title{Prescribed equality in off-grid extensions of G449}
\usepackage{etoolbox}
\makeatletter
\providecommand{\subtitle}[1]{% add subtitle to \maketitle
  \apptocmd{\@title}{\par {\large #1 \par}}{}{}
}
\makeatother
\subtitle{Up to six new points, with a marked nonadjacent pair sharing a
color}
\author{Aafi Khan Malek (Independent Researcher, Ahmedabad, India)}
\date{November 2026}

\begin{document}
\maketitle

\section{Equality can be prescribed with up to six off-grid
additions}\label{equality-can-be-prescribed-with-up-to-six-off-grid-additions}

This is a written lemma for the pinned G449: let x be an original vertex
and z a new off-grid point not at unit distance from x. Adjoin z and up
to five other distinct off-grid points. The full induced unit-distance
graph has a proper five-coloring with x and z the same color. Here the
grid is \(L=\tfrac17\mathbf{Z}[\zeta_{42}]\) in the exact complex
coordinates of the pinned G449, with all 449 original sites in L;
``off-grid'' means outside L. In particular no distance-two
original-to-marked terminal gadget in this off-grid subclass forces the
terminal colors apart. This is stronger in this subclass than ordinary
five-colorability, not an improvement to the whole-plane bound.

\subsection{Proof}\label{proof}

The previous off-grid degree lemma for this exact G449 says every
off-grid site has at most two original neighbors. Choose a checked base
coloring such that no original neighbor of z has the color of x, using
the explicit pair and simultaneous triple palette witnesses (or any base
palette if z has none). Give z the color of x. No neighbor of z was x
because their distance is not one.

Consider each remaining point's list of available colors after the base
and z are fixed. A point near z, meaning at unit distance from it, loses
at most three colors (two original neighbors and z), so it has a list of
size at least two. A far point loses at most two, so its list has size
at least three. It suffices to show that the induced graph on the at
most five remaining points is colorable from any lists satisfying these
respective bounds.

Suppose otherwise and select an inclusion-minimal uncolorable induced
subset S. Every near point in S has degree at least two, and every far
point degree at least three in S. Indeed, if it had fewer neighbors than
available colors, color the smaller induced graph and extend greedily.

Two elementary geometric facts are used:

\begin{enumerate}
\def\labelenumi{\arabic{enumi}.}
\tightlist
\item
  Distinct points have at most two common unit neighbors, as their unit
  circles have at most two intersection points. The unit-distance graph
  contains no K4: for two unit-separated points, their two common unit
  neighbors are separated by \(\sqrt{3}\), not 1.
\item
  On the unit circle centered at z, the only given point's unit
  neighbors have arguments shifted by \(\pm60^\circ\). The unit chord
  graph is therefore a disjoint union of subgraphs of six-cycles. Any
  set of at most five points on this circle induces a forest.
\end{enumerate}

\subsubsection{Up to four remaining
points}\label{up-to-four-remaining-points}

If \(|S|\le3\), the minimum-degree bounds exclude any far point; the
near-only induced subgraph is a forest, also contradicting minimum
degree two. Thus \(|S|=4\). Each far point is then universal in S. With
at least three far points, S induces K4, impossible. With exactly two
far points, both near points have the two far points as common unit
neighbors, in addition to z, impossible. With exactly one far point, the
other three points each need a near neighbor as well as this universal
far neighbor. Their induced forest must then be a path of length two.
The two endpoints already share z and the path's middle as common unit
neighbors, so the far point cannot be adjacent to both. Finally, with no
far point, the induced graph is a forest and fails the
minimum-degree-two condition. This rules out \(|S|\le4\).

\subsubsection{Five remaining points}\label{five-remaining-points}

From now on \(|S|=5\). Two further consequences of the elementary facts
are useful. A far point u that has two \emph{adjacent} near neighbors
a,b cannot have any third near neighbor: z and u are the two common unit
neighbors of a,b, and hence a,b are the two common unit neighbors of
z,u. Also, a near path a--b--c has endpoints a,c with common unit
neighbors z,b, so no far point can neighbor both endpoints.

\textbf{Five far.} Minimum degree three on five vertices means the
missing edges form a matching. If it has fewer than two edges, there is
a K4. If it has exactly two edges, their endpoints induce a four-cycle
and the remaining vertex is universal. Such a cycle cannot lie on the
unit circle about the universal vertex.

\textbf{Four far.} The four far points induce a graph of minimum degree
at least two. It is either a four-cycle, a diamond (K4 minus one edge),
or K4: the missing edges form a matching. K4 is excluded. In the
four-cycle case every far point needs its third neighbor at the unique
near point, giving a four-cycle on its unit circle. In the diamond case
the two degree-two far points must both neighbor the near point, but
already share the other two far points as common unit neighbors. Both
cases are impossible.

\textbf{Three far.} Each of the three far points has at least one far
neighbor, as there are only two near points. If the far subgraph is a
path, its endpoints must both neighbor both near points to reach degree
three. The two near points would then have those two endpoints and z as
common unit neighbors. Hence the far points form a unit equilateral
triangle. Each far point must see at least one near point, and no near
point can see all three (K4).

If both near points see two far points, they cannot see the same pair
(the third far point is already a common neighbor of that pair). Their
far pairs must then be distinct, with one common far point p.~Write the
other two far points as q,r. The near points are the opposite
equilateral apices across the pairs p,q and p,r: in complex coordinates
the near points are \(p+q-r\) and \(p+r-q\). Their separation is exactly
two and their midpoint is p.~Thus their unique common unit neighbor is
p, not z. If the near points are nonadjacent, each needed at least two
far neighbors, so this case is finished. If the near points are
adjacent, each needs at least one far neighbor; the only remaining
possibility to cover all far points has disjoint far-neighbor sets of
sizes one and two. Name the far triangle u,v,w, with one near point a
seeing u alone and the other b seeing v,w. The location of b is
\(v+w-u\), the other apex over the edge v,w. Thus v,w are already the
two common unit neighbors of u,b. But a needs to be another, since it is
adjacent to both u and b. This is impossible.

\textbf{Two far.} If the three near points have no near edge, each must
see both far points, contradicting at most two common unit neighbors of
the far pair. In either remaining near subgraph there is a near edge,
and the far pair must be adjacent: if not, each far point would need to
see all three near points, including an adjacent pair, which was ruled
out above.

If the near subgraph is a two-edge path a--b--c, each far point must see
at least two near points, but neither can see both endpoints a,c.~Since
each endpoint needs a far neighbor, the far points see a,b and b,c
respectively. Their positions are forced as the other equilateral apices
across these two edges. Put z at (0,0), b at (1,0), and a,c at
\((1/2,\pm\sqrt{3}/2)\). The other apices are \((3/2,\pm\sqrt{3}/2)\),
separated by \(\sqrt{3}\). They are not adjacent, a contradiction.

If the near subgraph is a single edge a--b and an isolated point c, then
c must see both far points. Each far point must see one of a,b as well,
and cannot see both, since they are an adjacent near pair and it already
sees c.~The points a,b both need a far neighbor. Label the far points
u,v so that u sees a,c and v sees b,c.~The induced graph is exactly the
four-cycle a--b--v--u--a plus c adjacent to u,v. Choose any color for
c.~Now each of u,v still has at least two available colors, and a,b had
at least two. An even cycle is colorable from arbitrary lists of size at
least two: shrink to size two; if all lists coincide, alternate, and
otherwise start with a color missing from the list of a chosen adjacent
last vertex, then color in cyclic order greedily. This colors the
putatively uncolorable subset.

\textbf{One far.} Each of the four near points needs a near neighbor.
Their induced forest is then either a four-vertex path or a pair of
disjoint edges. Any three near points contain an adjacent pair. The far
point has degree at least three, but it cannot have three near neighbors
including an adjacent pair.

\textbf{Zero far.} The near graph is a forest, impossible with minimum
degree two.

This exhausts all sizes and gives the required extension.

The proof assumes the archived off-grid degree result and the archived
exact palette witnesses, not an unstated plane theorem. No new
computation, solver, or Lean proof was used.

\section{Provenance and scope}\label{provenance-and-scope}

Correspondence:
\href{mailto:aafimalek2032@gmail.com}{\nolinkurl{aafimalek2032@gmail.com}}.
The research and writing in this note were developed with AI assistance.

The exact off-grid degree and simultaneous-palette prerequisites are
documented in the
\href{https://github.com/Aafimalek/novel-discoveries/tree/487e4dfc1f1bcbd536a4d03c814db4d0d65da779/hadwiger-nelson/verification/g449/prime-seven-catalogue/one-point-distance2-terminals}{one-point
terminal theorem}. The
\href{https://github.com/Aafimalek/novel-discoveries/tree/294c4596476e933f75703479fab2ea0f87e3723e/hadwiger-nelson/verification/g449/prime-seven-catalogue/six-point-extension-theorem}{ordinary
six-point extension theorem} has no off-grid restriction but does not
prescribe equal colors. This written note does not change the
whole-plane chromatic bounds and is not formalized in Lean.

\end{document}
